% Draft IEEEtran skeleton — compile with: pdflatex + bibtex
% Requires IEEEtran.cls from CTAN / IEEE
\documentclass[journal]{IEEEtran}
\usepackage{amsmath,amssymb,graphicx,cite,url,booktabs}
\begin{document}

\title{RECYM: An Automated Multi-Feeder Workflow for Demand Allocation and Load-Flow Studies in CYMDIST Distribution Networks}

\author{
  \IEEEauthorblockN{Mac Tapia Ccosyo\IEEEauthorrefmark{1}}
  \IEEEauthorblockA{\IEEEauthorrefmark{1}Electro Dunas S.A.A., Peru \\
  Email: mtapia@electrodunas.com}
}

\maketitle

\begin{abstract}
Distribution utilities need repeatable procedures to clean network models, allocate measured substation demand, inject important-customer and prospective spot loads, and compare situational versus projected operating cases. This paper presents RECYM, a multi-feeder software suite that couples CYMDIST/CymPy with a seven-step SPA and batch pipeline. On feeder PA217 (22.9\,kV, Electro Dunas), the campaign closed with zero diagnostic problems, eleven important customers verified within tight tolerances, headroom-balanced allocation ($\Delta\approx0\%$), a 1420\,kW prospective spot load, and successful dual load flows.
\end{abstract}

\begin{IEEEkeywords}
Distribution networks, load allocation, load flow analysis, CYMDIST, power system modeling, software tools.
\end{IEEEkeywords}

\section{Introduction}
% Full text: see MANUSCRIPT.md
Modern distribution planning relies on detailed feeder models~\cite{Kersting2017,Gonen2014}.
Commercial packages such as CYMDIST provide load allocation and unbalanced load-flow engines~\cite{CYME92}.
This paper presents RECYM and validates it on Electro Dunas feeder PA217.

\section{Related Work}
See~\cite{Dugan2011OpenDSS,Thurner2018Pandapower,Peppanen2015AMI,Rylander2015Hosting,IEEE1547_2018}.

\section{System Architecture}
Layered SPA/API/pipeline/CymPy design; multi-feeder JSON configuration; COM load allocation/flow~\cite{CYME92}.

\section{Proposed Workflow}
Seven steps: context, quality gate, customers+allocation, spot load, dual LF, reports, suite.

\section{Case Study}
PA217, $V_{LL}=22.9$\,kV, headroom $P=9537.88$\,kW, prospective load 1420\,kW.

\section{Results and Discussion}
Zero diagnostic messages; 11/11 precision OK; allocation balance $\Delta\approx0\%$; situational and projected LF OK; 21/21 closure steps.

\section{Conclusion}
RECYM automates CYMDIST demand campaigns with auditable artifacts and dual-scenario load flows.

\appendices
\section{Batch Sequence}
validate\_inputs $\rightarrow$ \ldots $\rightarrow$ demand\_allocation $\rightarrow$ report.

\bibliographystyle{IEEEtran}
\bibliography{references}

\end{document}
